Sure, I can help you with that! Here's an example code using TikZ to draw a shape with overlapping regions and adjustable opacity:
```
\documentclass{article}
\usepackage{tikz}
\begin{document}
\begin{tikzpicture}
% Define variables for shape dimensions
\def\widthA{3cm}
\def\heightA{2cm}
\def\widthB{4cm}
\def\heightB{3cm}
% Draw first rectangle
\fill[blue,opacity=0.5] (0,0) rectangle (\widthA,\heightA);
% Draw second rectangle
\fill[red,opacity=0.5] (\widthA-\widthB,0) rectangle (\widthA,\heightB);
% Draw third rectangle
\fill[yellow,opacity=0.5] (0,\heightA-\heightB) rectangle (\widthB,\heightA);
% Draw fourth rectangle
\fill[green,opacity=0.5] (\widthA-\widthB,\heightA-\heightB) rectangle (\widthA,\heightB);
\end{tikzpicture}
\end{document}
```
In this code, we define four rectangles with different colors and opacities. The first rectangle is blue and has a width and height of `\widthA` and `\heightA`, respectively. The second rectangle is red and overlaps with the first one on the right side. We adjust its position by subtracting `\widthB` from `\widthA`. The third rectangle is yellow and overlaps with the first one on the bottom side. We adjust its position by subtracting `\heightB` from `\heightA`. Finally, the fourth rectangle is green and overlaps with both the first and second rectangles.
You can modify the values of `\widthA`, `\heightA`, `\widthB`, and `\heightB` to change the dimensions of the rectangles. You can also adjust the `opacity` value for each rectangle to control the transparency level.